\documentclass[11pt,a4paper]{article}

\usepackage[spanish,es-noshorthands]{babel}
\usepackage[utf8]{inputenc}
\usepackage[T1]{fontenc}
\usepackage[margin=2.5cm]{geometry}
\usepackage{longtable}
\usepackage{booktabs}
\usepackage{array}
\usepackage{xcolor}
\usepackage{hyperref}
\usepackage{titlesec}
\usepackage{enumitem}
\usepackage{fancyhdr}
\usepackage{parskip}

\definecolor{pantone186}{RGB}{200,16,46}

\hypersetup{
  colorlinks=true,
  linkcolor=pantone186,
  urlcolor=pantone186,
  citecolor=pantone186,
  pdftitle={ERS - Sistema de Mesa de Ayuda (Ticketera)},
  pdfauthor={Equipo Ticketera Nestjs}
}

\titleformat{\section}{\normalfont\Large\bfseries\color{pantone186}}{\thesection}{1em}{}
\titleformat{\subsection}{\normalfont\large\bfseries}{\thesubsection}{1em}{}

\pagestyle{fancy}
\setlength{\headheight}{14pt}
\fancyhf{}
\renewcommand{\headrulewidth}{0.4pt}
\fancyhead[L]{ERS — Sistema de Mesa de Ayuda (Ticketera)}
\fancyhead[R]{\thepage}
\fancyfoot[C]{Documento interno — confidencial}

\newcolumntype{L}[1]{>{\raggedright\arraybackslash}p{#1}}

\newcommand{\reqtable}[1]{%
  \begin{longtable}{@{}>{\bfseries}p{2.1cm} L{2.6cm} L{7.7cm} p{1.4cm}@{}}
  \toprule
  \textbf{ID} & \textbf{Categoría} & \textbf{Descripción} & \textbf{Prioridad} \\
  \midrule
  \endhead
  #1
  \bottomrule
  \end{longtable}
}

\title{\color{pantone186}\textbf{Especificación de Requisitos de Software (ERS)}\\[0.3em]
\Large Sistema de Mesa de Ayuda (Ticketera)}
\author{Proyecto Ticketera Nestjs}
\date{22 de septiembre de 2026 \\ Versión 1.0}

\begin{document}

\maketitle
\thispagestyle{empty}

\begin{center}
\begin{tabular}{ll}
\textbf{Documento:} & Especificación de Requisitos de Software (ERS) \\
\textbf{Norma de referencia:} & IEEE 830-1998 (estructura adaptada) \\
\textbf{Estado:} & Borrador para revisión \\
\textbf{Alcance del documento:} & Backend completo del sistema (5 microservicios) \\
\end{tabular}
\end{center}

\newpage
\tableofcontents
\newpage

% =====================================================================
\section{Introducción}

\subsection{Propósito}
Este documento especifica los requisitos funcionales y no funcionales del
backend del sistema de mesa de ayuda (ticketera) del proyecto
\textbf{Ticketera Nestjs}. Su objetivo es servir de referencia única y
verificable para el equipo de desarrollo, control de calidad, y las
revisiones de seguridad, infraestructura y mantenibilidad que ya se han
realizado sobre el sistema. Los requisitos aquí descritos reflejan el
comportamiento real e implementado del sistema, documentado a partir de la
arquitectura, el código y los hallazgos de auditoría registrados en el
repositorio.

\subsection{Alcance del documento}
La especificación cubre exclusivamente el \textbf{backend}: los cinco
procesos que componen el sistema (\texttt{api-gateway}, \texttt{bff-web},
\texttt{users-service}, \texttt{tickets-service},
\texttt{notifications-service}) y el servicio de ingesta de correo
(\texttt{email-ingestion-service}). No especifica el frontend, dado que el
proyecto no incluye un cliente propio en el repositorio (solo mockups de
referencia pendientes de aprobación para la vista del solicitante).

\subsection{Definiciones, acrónimos y abreviaturas}

\begin{longtable}{@{}>{\bfseries}p{3.2cm} L{11cm}@{}}
\toprule
\textbf{Término} & \textbf{Definición} \\
\midrule
\endhead
ERS & Especificación de Requisitos de Software \\
RF & Requisito Funcional \\
RNF & Requisito No Funcional \\
SLA & Service Level Agreement — plazo de atención comprometido para un ticket \\
BFF & Backend For Frontend — servicio de composición de pantallas \\
RPC & Remote Procedure Call — invocación entre microservicios vía TCP \\
JWT & JSON Web Token — mecanismo de autenticación del gateway \\
ARCO & Derechos de Acceso, Rectificación, Cancelación y Oposición (Ley 21.719, Chile) \\
CRUD & Create, Read, Update, Delete \\
DTO & Data Transfer Object \\
UoW & Unit of Work — patrón de manejo transaccional \\
GIN/trgm & Índices de PostgreSQL para búsqueda de texto por similitud \\
ADMIN, SUPERVISOR, AGENTE, REQUESTER & Roles de usuario del sistema, de mayor a menor privilegio \\
\bottomrule
\end{longtable}

\subsection{Referencias}
\begin{itemize}[leftmargin=1.2cm]
  \item \texttt{CLAUDE.md} — guía de arquitectura y decisiones del repositorio.
  \item \texttt{libs/common/src/contracts/ticket-permissions.ts} — reglas de autorización por rol y estado.
  \item \texttt{docs/proceso/} — planificación, desarrollo, implementación y marcha blanca.
  \item Ley N.º 21.719 de Protección de Datos Personales (Chile).
  \item IEEE Std 830-1998, \emph{Recommended Practice for Software Requirements Specifications}.
\end{itemize}

% =====================================================================
\section{Descripción general}

\subsection{Perspectiva del producto}
El sistema es un backend de mesa de ayuda construido como un \textbf{monorepo
NestJS} con cinco procesos independientes que se comunican por transporte TCP
de \texttt{@nestjs/microservices}, respaldados por una única base de datos
PostgreSQL. El \texttt{api-gateway} es el único proceso con puerto HTTP
expuesto; el resto vive en una red interna sin acceso directo desde fuera.

\begin{center}
\begin{tabular}{ll}
\textbf{api-gateway} & Borde HTTP: CORS, helmet, rate limit, JWT, Swagger \\
\textbf{bff-web} & Composición de pantallas para lecturas, sin base de datos propia \\
\textbf{users-service} & Identidad, roles, derechos ARCO \\
\textbf{tickets-service} & Dominio de tickets: creación, asignación, SLA, estados \\
\textbf{notifications-service} & Observador de eventos de dominio, envío de avisos \\
\textbf{email-ingestion-service} & Conversión de correos entrantes en tickets \\
\end{tabular}
\end{center}

\subsection{Funciones del producto}
En términos generales, el sistema permite: autenticar usuarios; crear
tickets manualmente o por correo; clasificar su urgencia automáticamente;
calcularles un SLA; asignarlos a un agente según distintas estrategias;
registrar comentarios y cambios de estado con una máquina de estados
gobernada por rol; notificar a los interesados ante cada cambio; y ejercer
los derechos ARCO sobre los datos personales almacenados.

\subsection{Características de los usuarios}

\begin{longtable}{@{}>{\bfseries}p{2.6cm} L{11.6cm}@{}}
\toprule
\textbf{Rol} & \textbf{Perfil y permisos} \\
\midrule
\endhead
REQUESTER & Persona que reporta incidencias. Ve y actúa solo sobre sus propios tickets, con transiciones que dependen del estado de origen. \\
AGENTE & Resuelve tickets asignados o sin asignar de su especialidad. Comenta y cambia estado dentro de las reglas de autorización. \\
SUPERVISOR & Visibilidad ampliada: reportes, reasignación, estado de salud del sistema. \\
ADMIN & Gestión de usuarios (alta, roles, capacidad), configuración y todo lo anterior. \\
\bottomrule
\end{longtable}

\subsection{Restricciones generales}
\begin{itemize}[leftmargin=1.2cm]
  \item Toda la interfaz de usuario queda fuera de este backend; no existe cliente propio en el repositorio.
  \item El sistema debe ejecutarse sobre PostgreSQL 16 con \texttt{synchronize} deshabilitado: todo cambio de esquema requiere una migración escrita a mano.
  \item El transporte entre procesos es TCP interno sin broker de mensajería, por lo que los eventos de integración no son durables.
  \item El envío saliente real de correos electrónicos no está conectado a un servidor SMTP.
\end{itemize}

\subsection{Supuestos y dependencias}
\begin{itemize}[leftmargin=1.2cm]
  \item Se asume una única instancia por servicio en producción (el scheduler de SLA no soporta réplicas sin un lock distribuido adicional).
  \item Se asume disponibilidad de un buzón IMAP dedicado (\texttt{soporte@cmm.uchile.cl}) para la ingesta de correo, con contraseña de aplicación de Google Workspace.
  \item Se asume que el primer usuario ADMIN se crea por consola antes de exponer el sistema, dado que la API exige un ADMIN autenticado para crear usuarios.
\end{itemize}

% =====================================================================
\section{Requisitos específicos}

\subsection{Requisitos funcionales}

Los requisitos funcionales se numeran de \textbf{RF-01} a \textbf{RF-24} y se
agrupan por área funcional. Prioridad: \textbf{A} (alta / implementada y
crítica), \textbf{M} (media).

\reqtable{
RF-01 & Autenticación & El sistema debe permitir iniciar sesión con correo y contraseña, devolviendo un token de acceso (JWT, 15 min) y un token de refresco (7 días). & A \\
RF-02 & Autenticación & El sistema debe permitir renovar la sesión mediante el token de refresco sin exigir credenciales nuevamente. & A \\
RF-03 & Usuarios & El sistema debe permitir a un ADMIN crear usuarios con rol, habilidades (\emph{skills}) y capacidad máxima de tickets concurrentes. & A \\
RF-04 & Usuarios & El sistema debe permitir crear el primer usuario ADMIN mediante un comando de consola, sin depender de la API. & A \\
RF-05 & Usuarios & El sistema debe permitir a un ADMIN consultar, actualizar y dar de baja (desactivar) usuarios existentes. & M \\
RF-06 & Tickets & El sistema debe permitir a un usuario autenticado crear un ticket con título, descripción y categoría. & A \\
RF-07 & Tickets & El sistema debe clasificar automáticamente la urgencia del ticket analizando el texto del reporte mediante una heurística de señales de urgencia. & A \\
RF-08 & Tickets & El sistema debe calcular un SLA (plazo de atención) para cada ticket según su prioridad y categoría, aplicando una estrategia de cálculo. & A \\
RF-09 & Tickets & El sistema debe asignar automáticamente el ticket a un agente disponible según una estrategia configurable (por habilidad, por menor carga o round-robin). & A \\
RF-10 & Tickets & El sistema debe permitir a un SUPERVISOR reasignar manualmente un ticket a otro agente. & M \\
RF-11 & Tickets & La asignación de un ticket al agente que ya lo tiene asignado debe ser una operación idempotente (no generar un nuevo aviso). & A \\
RF-12 & Tickets & El sistema debe permitir cambiar el estado de un ticket únicamente entre las transiciones permitidas por rol y por el estado de origen actual. & A \\
RF-13 & Tickets & El sistema debe permitir agregar comentarios a un ticket, autorizando la acción según el estado real del ticket y el rol del autor. & A \\
RF-14 & Tickets & El sistema debe registrar un historial inmutable de cada cambio de estado del ticket (quién, cuándo, de qué estado a cuál). & A \\
RF-15 & Tickets & El sistema debe permitir buscar y filtrar tickets por estado, vencimiento de SLA, texto libre y rango de fechas de creación. & M \\
RF-16 & Tickets & Los listados de tickets deben soportar paginación y ordenamiento por columnas explícitamente permitidas. & M \\
RF-17 & Notificaciones & El sistema debe notificar al solicitante y al agente asignado (según corresponda) ante cada cambio de estado o comentario nuevo, identificando al interesado real y no a quien ejecutó la acción. & A \\
RF-18 & Notificaciones & El sistema debe degradar el envío de notificaciones a canal in-app cuando el servicio de resolución de destinatarios no esté disponible. & M \\
RF-19 & Ingesta de correo & El sistema debe convertir automáticamente un correo entrante en un ticket nuevo, o en un comentario si el correo es una respuesta a un ticket existente. & A \\
RF-20 & Ingesta de correo & El sistema debe descartar correos que no representen un pedido legítimo de una persona (bucles, remitentes automáticos, dominios no autorizados). & A \\
RF-21 & Ingesta de correo & El sistema debe dar de alta automáticamente, con rol REQUESTER, a un remitente desconocido cuyo dominio esté autorizado. & M \\
RF-22 & Monitoreo & El sistema debe barrer periódicamente los tickets abiertos para detectar incumplimientos de SLA y marcarlos. & A \\
RF-23 & Monitoreo & El sistema debe exponer el estado de los circuitos (dependencias caídas o degradadas) a usuarios con rol SUPERVISOR o ADMIN. & M \\
RF-24 & Privacidad (ARCO) & El sistema debe permitir a un usuario ejercer los derechos de acceso, rectificación, cancelación, oposición y portabilidad sobre sus datos personales, dejando registro auditable de cada solicitud. & A \\
}

\subsection{Requisitos no funcionales}

Los requisitos no funcionales se numeran de \textbf{RNF-01} a
\textbf{RNF-24}, agrupados por atributo de calidad.

\reqtable{
RNF-01 & Seguridad & Toda decisión de autorización que dependa del dato (por ejemplo, si un ticket pertenece al solicitante) debe resolverse en el microservicio dueño del dato, nunca inferirse en el borde. & A \\
RNF-02 & Seguridad & Toda función de autorización debe enumerar explícitamente los casos permitidos y lanzar una excepción en cualquier otro caso; no debe existir una rama final que permita por omisión. & A \\
RNF-03 & Seguridad & Toda comunicación RPC entre microservicios debe ir autenticada con un secreto compartido (\texttt{RPC\_SHARED\_SECRET}), distinto del secreto de firma de JWT. & A \\
RNF-04 & Seguridad & Los microservicios internos no deben exponer su puerto TCP fuera de la red local del proceso (\texttt{MICROSERVICE\_BIND\_HOST}). & A \\
RNF-05 & Seguridad & El gateway debe ser el único proceso con acceso a internet y el único que conoce el secreto de firma del JWT; no debe tener acceso directo a la base de datos. & A \\
RNF-06 & Seguridad & Las contraseñas de usuario deben almacenarse con hashing, nunca en texto plano. & A \\
RNF-07 & Seguridad & La documentación interactiva de la API (Swagger) debe poder deshabilitarse por configuración y, en producción, exigir autenticación básica con usuario y contraseña obligatorios. & A \\
RNF-08 & Seguridad & Las búsquedas de texto libre deben tratar el término de búsqueda como parámetro, nunca interpolado en la consulta SQL. & A \\
RNF-09 & Disponibilidad & Toda llamada saliente entre procesos debe protegerse con un circuit breaker que solo cuente fallos de infraestructura contra su umbral, nunca respuestas 4xx legítimas. & A \\
RNF-10 & Disponibilidad & Un error de autorización o de negocio devuelto por un servicio dependiente no debe transformarse en un error 500 al cruzar otro proceso. & A \\
RNF-11 & Disponibilidad & Los cinco procesos y el servicio de ingesta deben exponer un chequeo de vida (\emph{liveness}) y uno de disponibilidad (\emph{readiness}) independientes. & A \\
RNF-12 & Rendimiento & Los tiempos de espera de lock, de sentencia SQL y del circuit breaker deben estar escalonados de menor a mayor, de forma que ninguna capa espere más que la que la contiene. & A \\
RNF-13 & Rendimiento & Ninguna transacción de base de datos debe permanecer abierta mientras se realiza una llamada de red a otro proceso. & A \\
RNF-14 & Rendimiento & Las consultas por atributos de uso frecuente deben poder resolverse usando un índice, verificado con \texttt{EXPLAIN} y \texttt{enable\_seqscan=off}. & M \\
RNF-15 & Rendimiento & La resolución de destinatarios de un lote de notificaciones debe realizarse en una única llamada RPC, no una por destinatario. & M \\
RNF-16 & Escalabilidad & El sistema debe documentar explícitamente sus límites de escalado horizontal (resolución fija de \texttt{host:port} en los clientes RPC, límite de \emph{throttling} por instancia). & M \\
RNF-17 & Mantenibilidad & El dominio de cada microservicio debe organizarse en las capas \texttt{domain/}, \texttt{application/} e \texttt{infrastructure/}, sin que el dominio dependa de \texttt{@nestjs/microservices} ni de HTTP. & A \\
RNF-18 & Mantenibilidad & Las reglas de autorización de tickets deben existir en un único lugar del código, compartido entre el backend y el BFF, para que la interfaz y el caso de uso nunca se desincronicen. & A \\
RNF-19 & Confiabilidad & Debe existir un mecanismo de respaldo de la base de datos verificado y con retención, y un procedimiento de restauración probado que cuente filas por tabla. & A \\
RNF-20 & Confiabilidad & Las migraciones de esquema deben aplicarse mediante un servicio dedicado antes de que arranque cualquier proceso que dependa de la base. & A \\
RNF-21 & Portabilidad & La configuración de cada proceso debe provenir de variables de entorno validadas contra un esquema explícito; una variable no declarada debe impedir el arranque. & A \\
RNF-22 & Portabilidad & El sistema completo debe poder desplegarse mediante contenedores Docker orquestados por \texttt{docker compose}, con cada servicio recibiendo solo los secretos que efectivamente usa. & A \\
RNF-23 & Usabilidad & Todo error de validación devuelto al cliente debe indicar qué campo específico es inválido y por qué, sin exponer detalles internos del sistema. & M \\
RNF-24 & Cumplimiento normativo & El tratamiento de datos personales debe ajustarse a la Ley 21.719 (Chile): bloqueo temporal de tratamiento, portabilidad y trazabilidad auditable de cada derecho ejercido. & A \\
}

% =====================================================================
\section{Alcances}

El presente sistema, en su versión especificada por este documento,
comprende:

\begin{itemize}[leftmargin=1.2cm]
  \item El backend completo de la mesa de ayuda: los cinco microservicios
    (\texttt{api-gateway}, \texttt{bff-web}, \texttt{users-service},
    \texttt{tickets-service}, \texttt{notifications-service}) y el servicio
    de ingesta de correo (\texttt{email-ingestion-service}).
  \item El ciclo de vida completo de un ticket: creación (manual o por
    correo), clasificación de urgencia, cálculo de SLA, asignación,
    comentarios, cambios de estado y notificación a los interesados.
  \item La gestión de usuarios y roles, incluyendo el ejercicio de los
    derechos ARCO y portabilidad de datos personales según la Ley 21.719.
  \item Los mecanismos transversales de seguridad, resiliencia
    (\emph{circuit breaker}), observabilidad y auditoría descritos en los
    requisitos no funcionales.
  \item La infraestructura de despliegue mediante contenedores Docker,
    incluyendo el servicio de migración de base de datos y los scripts de
    respaldo y restauración.
  \item Mockups de referencia (no funcionales) para las pantallas de
    agente/supervisor: login, listado de tickets, detalle de ticket y
    reportes, en versión web y mobile.
\end{itemize}

% =====================================================================
\section{Limitaciones}

Las siguientes limitaciones son decisiones de diseño ya tomadas y
documentadas en el proyecto, no omisiones accidentales. Se listan junto con
su costo asumido para que no se re-discutan sin recordar por qué existen.

\begin{itemize}[leftmargin=1.2cm]
  \item \textbf{No incluye un frontend propio.} El repositorio no contiene
    cliente web ni mobile; solo mockups de referencia. La vista del
    solicitante (rol REQUESTER) está deliberadamente sin diseñar, pendiente
    de aprobación de los \emph{stakeholders}.
  \item \textbf{Una sola base de datos PostgreSQL} para los cinco procesos,
    en vez de una por servicio. No es posible escalar o migrar los datos de
    un servicio de forma independiente del resto.
  \item \textbf{Sin sistema de mensajería intermedio (broker).} Los eventos
    de integración no son durables: si \texttt{notifications-service} está
    caído en el instante exacto en que ocurre un evento, esa notificación
    puntual se pierde (el cambio en el ticket queda registrado igual).
  \item \textbf{El envío real de correo saliente no está implementado.} El
    sistema determina qué correo enviar y a quién, pero el paso final hacia
    un servidor SMTP queda pendiente de conexión.
  \item \textbf{Sin revocación inmediata de sesión.} El token de acceso es
    la fuente de verdad durante sus 15 minutos de vida; desactivar un
    usuario no invalida una sesión ya emitida hasta que expira.
  \item \textbf{El \emph{scheduler} de SLA asume una única instancia por
    proceso.} Con réplicas activas del mismo servicio, dos procesos
    barrerían los tickets en paralelo sin coordinación adicional.
  \item \textbf{El secreto RPC es único y estático}, no por par de
    servicios: autentica al llamador pero no lo identifica individualmente.
    Un servicio comprometido puede invocar a cualquier otro del grupo
    interno.
  \item \textbf{No hay \emph{WAL archiving}.} El punto de recuperación
    (RPO) efectivo es el intervalo entre respaldos manuales, no minutos.
  \item \textbf{Los respaldos se almacenan en el mismo disco que la base de
    datos}, sin copia remota automatizada.
  \item \textbf{El escalado horizontal de los servicios internos no reparte
    carga.} El cliente RPC de cada proceso resuelve \texttt{host:port} una
    única vez y mantiene un socket persistente; agregar réplicas sin un
    balanceador TCP delante no distribuye tráfico.
  \item \textbf{El límite de \emph{rate limiting} es por instancia}, no
    compartido entre réplicas del gateway.
  \item \textbf{La suite de pruebas end-to-end (\texttt{test:e2e}) está
    declarada pero no implementada}; la cobertura actual es de pruebas
    unitarias y de integración parcial.
  \item \textbf{Alcance de la auditoría de privacidad (ARCO):} un hallazgo
    de auditoría pendiente documenta que el bloqueo temporal de privacidad
    aún no es consultado por el resolver de asignación ni por el
    procesamiento de notificaciones — el campo se persiste correctamente
    pero todavía no detiene el tratamiento real de datos en esos dos
    puntos.
\end{itemize}

\end{document}
